% Written in Part B for independent audit and integration in Part D.
% Configuration lemma and cusp theorem: root/D actual-file audits passed.

\begin{lemma}[Proved: Gorenstein log canonical exceptional configurations]
\label{D:Gorenstein-lc-configurations}
Let $(S,s)$ be a normal Gorenstein log canonical surface singularity
over an algebraically closed field $k$, in arbitrary characteristic.
On its minimal regular resolution, precisely one of the following
possibilities occurs:
\begin{enumerate}
\item the exceptional fibre is empty, or it is a simple normal
crossings tree of smooth rational curves, each of self-intersection
$-2$, and every discrepancy is zero;
\item the exceptional fibre is one smooth elliptic curve, of
discrepancy $-1$;
\item the exceptional fibre is a nodal cycle of rational curves,
and every discrepancy is $-1$.  A single irreducible nodal rational
curve is allowed in this description.
\end{enumerate}
After blowing up the node in the last single-component case, the
cycle has smooth components.  On every further regular resolution,
the elliptic or cyclic core persists; point blowups only subdivide
edges of the cyclic core or attach rational trees.  Every component
of a cyclic core still has discrepancy $-1$.
\end{lemma}
\begin{proof}
We use existence of surface resolutions, contraction of exceptional
smooth rational $(-1)$-curves, and factorization of a proper
birational morphism between regular surfaces into point blowups.
These hold in arbitrary characteristic, including over excellent
strictly henselian local surface rings; see
\cite[Theorem 54.14.5, Lemmas 54.16.8 and 54.17.1;
tags 0BGP, 0C2M and 0C5R]{StacksDSurfaceResolution}.
The intersection matrix of the exceptional curves is negative
definite by \cite[section 10.1, Theorem 10.1]{Kollar2013D}.
The fibre is connected, since a proper birational morphism to a
normal surface has direct image of its structure sheaf equal to
the structure sheaf, and hence has connected fibres.

Write the minimal resolution as $f:Z\to S$ and
\[
 K_Z=f^*K_S+\sum_i a_i E_i,\qquad
 M=(E_i\cdot E_j),\qquad b_i=K_Z\cdot E_i.
\]
Adjunction gives
$b_i=2p_a(E_i)-2-E_i^2\geq0$.
Indeed $E_i^2<0$; if $p_a(E_i)=0$, then $E_i$ is smooth rational
and minimality excludes self-intersection $-1$.
All off-diagonal entries of $M$ are nonnegative.
If $a^+$ denotes the positive part of the vector $a$, and $a^-$
its nonpositive part, then a nonzero $a^+$ would give
\[
 0\leq (a^+)^{\mathsf T}b
 =(a^+)^{\mathsf T}Ma^++(a^+)^{\mathsf T}Ma^-<0,
\]
by negative definiteness and the disjoint supports of the two parts.
Thus every $a_i\leq0$.  Gorensteinness makes the $a_i$ integers,
and log canonicity gives $a_i\geq-1$.  Hence
$a_i\in\{-1,0\}$.

If $a_i=0$ for some $i$, then
\[
 0\leq b_i=\sum_{j\ne i}a_j(E_j\cdot E_i)\leq0.
\]
Every component meeting $E_i$ therefore also has discrepancy zero.
Connectedness shows that all discrepancies are zero.  Adjunction
then gives $p_a(E_i)=0$ and $E_i^2=-2$ for every $i$.
There can be no cycle in the intersection multigraph: the vector
equal to one on a shortest cycle and zero elsewhere would have
self-intersection at least zero.  The same calculation rules out
an intersection number at least two between two components.
A triple intersection would give a three-cycle and is also
impossible.  Thus the fibre is a simple normal crossings rational
tree, as asserted in the first case.

Otherwise all $a_i=-1$, and adjunction becomes
\begin{equation}\label{D:eq-Gorenstein-lc-adjunction}
 2p_a(E_i)-2=-\sum_{j\ne i}E_j\cdot E_i.
\end{equation}
If one component has arithmetic genus at least one, this equation
and connectedness say that it is the only component and has
arithmetic genus one.  Otherwise every component is smooth
rational and has total intersection number two with the remaining
components.

We check the possible singularities of the reduced divisor
$E=\bigcup_i E_i$, rather than assuming a characteristic-zero
classification.  We have $K_Z+E=f^*K_S$, so the pair $(Z,E)$ is
log canonical.  A point of multiplicity $m\geq3$ in $E$ is
impossible: blowing it up creates a divisor of discrepancy
$1-m\leq-2$.  At a singular point of multiplicity two, two
distinct tangent lines give an ordinary node.  If the tangent line
is repeated, the first blowup creates a divisor of discrepancy
$-1$, and the strict transform meets it at a single point with
total intersection multiplicity two.  If the strict transform
has multiplicity at least two there, the reduced total boundary
already has multiplicity at least three.  Otherwise the strict
transform is smooth and tangent to the exceptional curve;
blowing up this tangency gives a point at which three reduced
boundary branches, counted with multiplicity, meet.  In either
case a further blowup creates discrepancy at most $-2$.
This contradicts log canonicity.  This elementary blowup argument
works also in characteristics two and three.
Consequently $E$ is nodal.

An integral nodal curve of arithmetic genus one is either smooth
elliptic or a rational curve with one node.  In the remaining case,
equation~\eqref{D:eq-Gorenstein-lc-adjunction} says that the connected
dual multigraph is two-regular, so it is a cycle.  This includes
two rational curves meeting transversely at two distinct points.
Blowing up the node of an irreducible rational curve gives the
same two-component configuration.  Its new component has
discrepancy $1-2=-1$.

Finally, factorization into point blowups shows the description
on any dominating regular resolution.  A blowup at a node of the
cycle subdivides its edge and has discrepancy $1-1-1=-1$.
A blowup at a smooth point of the reduced fibre attaches a
rational component.  Further blowups of an attached tree still
leave the original core as the unique cyclic, respectively
positive-genus, part.  This proves the final assertion.
\end{proof}

\begin{theorem}[Proved: tame ramification above a Gorenstein cusp cycle]
\label{D:cusp-tame}
Let $(S,s)$ be a normal Gorenstein log canonical surface singularity
over an algebraically closed field $k$.  Suppose that its minimal
resolution has a rational nodal cycle as exceptional fibre.
For every finite separable quasi-\'etale cover $S'\to S$ and every
component $E_i$ of this cycle, every extension $F$ of the associated
divisorial valuation is tame:
\[
 e(F/E_i)\text{ is prime to the characteristic exponent of }k,
 \qquad \kappa(F)/\kappa(E_i)\text{ is separable}.
\]
The same assertion holds for the components inserted by blowing
up nodes of the cycle.  The assertion is valid in every
characteristic; no assumption on the degree of the cover is made.
Here a cover is taken with normal integral source, or is treated
componentwise, and it is \'etale on the punctured germ.  For a
quasi-\'etale cover of an algebraic surface, we obtain this germ
by shrinking away the other finitely many branch points.
\end{theorem}
\begin{proof}
We give the local and group-theoretic comparisons explicitly.
The conclusion descends from a Galois closure, so we may assume
that the cover is Galois.  The Galois closure of its restriction
to the punctured surface is finite \'etale, and normalization
over the excellent surface is finite and quasi-\'etale.

\par\smallskip\noindent\textbf{Canonical divisors and strict localization.}\quad
At every codimension-one point the original cover is \'etale.
Thus compatible rational canonical forms give
$K_{S'}=\pi^*K_S$ as Weil divisors.  Since $K_S$ is Cartier,
the canonical reflexive sheaf on $S'$ is invertible.
A normal surface is Cohen--Macaulay, so $S'$ is Gorenstein.
This proves the property of the covering surface; we do not
assume it merely from the hypothesis on the target.

For a pair of divisorial valuations $F/E$ in a finite separable
surface function-field extension, let $e=e(F/E)$ and let
$\delta_F$ be the order of the determinant of differentials.
The calculation in Proposition~\ref{D:lc-elliptic-tame} gives
\begin{equation}\label{D:eq-cusp-discrepancy}
 a(F,S')=e\,a(E,S)+\delta_F,
 \qquad \delta_F\geq e-1.
\end{equation}
Equality in the second inequality holds exactly when $e$ is
prime to the characteristic exponent and the residue extension is
separable.  For completeness, at the regular generic points
write $t=u z^e$ for uniformizers and choose a separating residue
parameter $x$ on the target.  Then
\[
 d(u z^e)\wedge dx
 =z^{e-1}(eu\,dz\wedge dx+z\,du\wedge dx).
\]
The coefficient at order $e-1$ is nonzero exactly under those
two conditions.  This includes inseparable residue extensions.
The restriction of a divisorial valuation through a finite
function-field extension is again divisorial.  Formula
\eqref{D:eq-cusp-discrepancy} consequently proves that $S'$ is
log canonical.

Replace $\mathcal O_{S,s}$ by its strict henselization $A$.
It is an excellent normal local ring with residue field $k$.
Choose one connected component of the punctured cover, and let
$B$ be the normalization of $A$ in its field $L$.
It is finite over $A$ by excellence.  A finite algebra over a
henselian local ring is a product of local rings; since $B$ is
a domain it is local, and its residue field is $k$.  Thus $B$
is strictly henselian.  The chosen component is Galois with
its decomposition group, which we now denote by $G$.
Normality, the canonical-divisor comparison, Gorensteinness and
log canonicity persist under this ind-\'etale localization.
Equivalently, all the finite-presentation constructions and the
regular models used to compute a fixed discrepancy descend to
an \'etale neighborhood of the original algebraic singularity.
We continue to write $S=\operatorname{Spec}A$ and
$S'=\operatorname{Spec}B$, and put
$U=S\setminus\{s\}$ and $U'=S'\setminus\{s'\}$.
The latter is a connected finite \'etale $G$-cover of $U$.

\par\smallskip\noindent\textbf{The nonzero cycle quotient.}\quad
Take a regular resolution $Y\to S$ whose reduced exceptional
fibre $E$ is a cycle of smooth rational curves.  If the minimal
cycle consists of a single nodal curve, blow up its node once.
This operation leaves each old discrepancy equal to $-1$ and
gives the new component discrepancy $-1$ as well.
For a proper scheme over a henselian local ring, restriction
to the closed fibre is an equivalence on finite \'etale covers;
see \cite[Lemma 58.9.1, tag 0A48]{StacksBCuspLocal}.
Nilpotents do not change this category.  Therefore
\[
 \pi_1(Y)\simeq\pi_1(E)\simeq\widehat{\mathbb Z}.
\]
The last identification follows directly by normalizing the
nodal curve: each normalized component is $\mathbb P^1_k$,
whose finite \'etale covers are trivial, and the gluing at the
nodes is exactly the data of a finite covering of the cycle
graph.  The same gluing description shows that a rational tree
has trivial fundamental group and that attaching rational
trees does not change the fundamental group of a core.

Since $Y$ is normal, the dense-open inclusion $U\hookrightarrow Y$
induces a surjection on fundamental groups: a connected finite
\'etale cover of $Y$ is normal and integral, and remains
connected on its nonempty open preimage of $U$.
The open subgroup $\pi_1(U')\subset\pi_1(U)$ therefore has
finite-index, and hence nonzero, image in
$\widehat{\mathbb Z}$.

\par\smallskip\noindent\textbf{The core on the covering resolution.}\quad
Let $Z_0\to S'$ be the minimal regular resolution, with a node
blown up if necessary to obtain smooth components in its
cyclic case.  Resolve the normalized graph of the rational
map $Z_0\dashrightarrow Y$.  We obtain a regular surface $Z$
with a proper birational morphism $Z\to Z_0$ and a morphism
$Z\to Y$.  It is unchanged over $U'$.  Surface factorization
shows that $Z\to Z_0$ is a sequence of point blowups with
centres in the exceptional fibre.  Let $E'$ be the reduced
exceptional fibre of $Z\to S'$.
Proper henselian comparison is functorial, so the composite
\[
 \pi_1(U')\longrightarrow\pi_1(Z)
 \simeq\pi_1(E')\longrightarrow\pi_1(Y)
 \simeq\pi_1(E)
\]
is the finite-index nonzero homomorphism of the preceding
paragraph.  In particular $\pi_1(E')\to\pi_1(E)$ is nonzero.

Apply Lemma~\ref{D:Gorenstein-lc-configurations} to $S'$.
If its minimal exceptional divisor is empty or a rational tree,
then $E'$ is a point or a rational tree and $\pi_1(E')=1$,
a contradiction.
If its minimal fibre is elliptic, $E'$ consists of that
elliptic core with rational trees attached.  The irreducible
core maps into one component of $E$, or into a point.
Its map on fundamental groups factors through
$\pi_1(\mathbb P^1_k)=1$.  The attached trees contribute
nothing, again a contradiction.
Thus the covering singularity is also a cusp, and $E'$ has
a rational cycle core $C'$ with rational trees attached.
Every component of $C'$ has discrepancy $-1$.

The image of $C'$ must dominate every component of $E$.
Indeed, otherwise its connected reduced image is contained
in a proper connected subcurve of the smooth-component
cycle $E$, hence in a rational chain, or is a point.
Such a subcurve has trivial fundamental group.  Since
$\pi_1(C')\to\pi_1(E')$ is an isomorphism, this would make
$\pi_1(E')\to\pi_1(E)$ zero, contrary to what we have proved.
Consequently, for each target component $E_i$ there is a
component $F_i$ of the core $C'$ mapping dominantly to it.
This is an extension of the corresponding divisorial
valuation, with discrepancy $a(F_i,S')=-1$.

\par\smallskip\noindent\textbf{All extensions of each valuation.}\quad
Normalize $Y$ in $L$.  The extensions of the discrete
valuation at $E_i$ are the prime divisors above it, and the
Galois group $G$ acts transitively on them.  The normal
resolution $Z$ factors through this normalization.
The discrepancy $a(F,S')$ is intrinsic to the valuation
and is invariant under deck transformations of $S'$:
its canonical Cartier class is the pullback of the one on
$S$.  Since one extension $F_i$ has discrepancy $-1$,
all extensions $F/E_i$ have discrepancy $-1$.
The target discrepancy is also $-1$, and
\eqref{D:eq-cusp-discrepancy} now gives
\[
 -1=-e(F/E_i)+\delta_F,
 \qquad \delta_F=e(F/E_i)-1.
\]
The equality criterion proves tameness, with separable
residue extension.  Starting with a target resolution
in which any prescribed cycle node has already been
blown up proves the additional assertion.

Finally, strict henselization has not weakened the
valuation conclusion for the original algebraic cover.
The exceptional curves and their residue function fields
over the closed point remain the same after this base
change, since its residue field is already $k$; the
extension ramification indices and residue extensions
are the ones computed in the corresponding local
components.  More explicitly, the induced maps of the
boundary discrete valuation rings are ind-\'etale with
unchanged residue fields, so these indices and residue
extensions are preserved.  Thus tameness for those components gives
the stated assertion before strict localization.
\end{proof}

\begin{remark}[Exact local scope]
The proof uses strict henselization to make the finite
normal covering ring local and to compare finite \'etale
covers of a proper resolution with its exceptional fibre.
It does not assert an equivalence of punctured algebraic,
strictly henselian and completed neighborhoods, and no
completion comparison is needed.  The Gorenstein
hypothesis is used twice: to make every covering surface
Gorenstein by canonical Cartier pullback, and to make
the minimal-resolution discrepancies integral.  We have
not replaced it by an arbitrary log canonical hypothesis.
\end{remark}

\begin{corollary}[Proved: finite omitted klt and Gorenstein log canonical points]
\label{D:Gorenstein-lc-boundary}
Let $X$ be a normal proper integral surface over an algebraically
closed field, and let $Y\subset X$ be a finite set of closed
points.  Suppose that each point of $Y$ is either log terminal
or Gorenstein and log canonical.  Then $\pi_1(X\setminus Y)$ is topologically
finitely generated.  The singularities retained in
$X\setminus Y$ are unrestricted.
\end{corollary}
\begin{proof}
Let $q:\widetilde X\to X$ be the minimal regular resolution,
and additionally blow up the node of any irreducible cusp
cycle over $Y$.  This is a proper birational morphism from
a regular proper surface.  Such a surface is projective;
equivalently, we may use the normal proper surface
projectivity criterion recalled in
\cite[introduction, pp.~1--2]{Schroer1999D}, for which the
non-$\mathbb Q$-factorial locus of a regular surface is empty.
The only prime divisors in
$\widetilde X\setminus q^{-1}(X\setminus Y)$ are the
exceptional components over points of $Y$.

Every finite \'etale cover of $X\setminus Y$ normalizes to
a finite quasi-\'etale cover of $X$, since $X$ is excellent
and $Y$ has codimension two.  At an elliptic or cusp point,
the preceding lemma identifies all components on the
chosen resolution with its elliptic or cyclic core.
Proposition~\ref{D:lc-elliptic-tame} and
Theorem~\ref{D:cusp-tame} show that the wild inertia image
at each such component is zero in $\pi_1(X\setminus Y)$.

Every remaining omitted point is log terminal, either by
the stated hypothesis or by the all-zero case of the
configuration lemma.  In that all-zero case, subsequent
point blowups have discrepancy greater than $-1$.
Its punctured strict-local \'etale fundamental group is
finite by \cite[Theorem 5.3, p.~739]{CarvajalRojasYasuda2025B}.
That primary theorem applies to surface germs in every
characteristic, including two and three.  The wild inertia
image at every exceptional component over such a point
factors through this finite local group, and is therefore
finite.  A regular omitted point has no exceptional divisor
on our chosen resolution and causes no additional condition.

All divisorial boundary wild images on the fixed projective
modification are thus finite.  There is no original
codimension-one boundary, so
$Q(X,Y)=\pi_1(X\setminus Y)$, and
Proposition~\ref{D:finite-exceptional-wild} applies.
\end{proof}
